
\documentclass{article}
\usepackage{colortbl}
\usepackage{makecell}
\usepackage{multirow}
\usepackage{supertabular}

\begin{document}

\newcounter{utterance}

\centering \large Interaction Transcript for game `clean\_up', experiment `2\_hard\_3obj\_nl', episode 0 with Qwen3.6{-}35B{-}A3B{-}FP8{-}without{-}reasoning.
\vspace{24pt}

{ \footnotesize  \setcounter{utterance}{1}
\setlength{\tabcolsep}{0pt}
\begin{supertabular}{c@{$\;$}|p{.15\linewidth}@{}p{.15\linewidth}p{.15\linewidth}p{.15\linewidth}p{.15\linewidth}p{.15\linewidth}}
    \# & $\;$A & \multicolumn{4}{c}{Game Master} & $\;\:$B\\
    \hline

    \theutterance \stepcounter{utterance}  
    & & \multicolumn{4}{p{0.6\linewidth}}{
        \cellcolor[rgb]{0.9,0.9,0.9}{
            \makecell[{{p{\linewidth}}}]{
                \texttt{\tiny{[P1$\langle$GM]}}
                \texttt{Ik ben uw spelleider, en u speelt een samenwerkingsspel met het volgende raster als speelbord:} \\
\\ 
\texttt{```} \\
\texttt{ 1234567} \\
\texttt{╔══╤════╗} \\
\texttt{║◌◌│C◌◌◌║ 1} \\
\texttt{╟──┤◌L┌─╢ 2} \\
\texttt{║◌◌├──┤◌║ 3} \\
\texttt{╟──┘◌◌├─╢ 4} \\
\texttt{║◌◌P◌◌│◌║ 5} \\
\texttt{╟─┐◌┌─┼─╢ 6} \\
\texttt{║◌│◌│◌│◌║ 7} \\
\texttt{╚═╧═╧═╧═╝} \\
\\ 
\texttt{```} \\
\\ 
\texttt{* De bovenrand toont x{-}coördinaten die naar rechts oplopen, en de rechterrand y{-}coördinaten die naar beneden oplopen.} \\
\texttt{* De volgende objecten zijn willekeurig op uw raster geplaatst: 'C', 'L', 'P'.} \\
\\ 
\\ 
\texttt{De andere speler ziet een variatie van het speelbord, waarbij de objecten op verschillende willekeurige locaties zijn geplaatst. U kunt het bord van de andere speler niet zien en zij kunnen het uwe niet zien.} \\
\\ 
\\ 
\texttt{**Doel:**} \\
\\ 
\texttt{Beide spelers moeten de objecten op hun respectievelijke achtergrond zo verplaatsen dat identieke objecten op dezelfde coördinaten terechtkomen. U moet communiceren met de andere speler om een gemeenschappelijke doelconfiguratie overeen te komen.} \\
\\ 
\\ 
\texttt{**Regels:**} \\
\\ 
\texttt{* In elke beurt kunt u precies één van de volgende twee commando's sturen:} \\
\texttt{ 1. `ZEG: <BERICHT>`: om een bericht (alles tot de volgende regeleinde) naar de andere speler te sturen. Ik zal het doorsturen naar uw partner.} \\
\texttt{2. `BEWEEG: <OBJECT>, (<X>, <Y>)`: om een object naar een nieuwe positie te verplaatsen, waarbij `<X>` de kolom is en `<Y>` de rij. Ik zal u laten weten of uw zet succesvol was of niet.} \\
\texttt{* Als u zich niet aan het formaat houdt, of meerdere commando's tegelijk verstuurt, moet ik u straffen.} \\
\texttt{* Als beide spelers meer dan 8 strafpunten verzamelen, verliezen jullie allebei het spel.} \\
\texttt{* Het is van vitaal belang dat u met de andere speler communiceert over uw doelstatus! De *enige* manier waarop u uw strategie aan de andere speler kunt overbrengen is door gebruik te maken van het `ZEG: <BERICHT>` commando!} \\
\\ 
\\ 
\texttt{**Objecten verplaatsen**} \\
\\ 
\texttt{* U kunt alleen objecten verplaatsen naar cellen binnen de grenzen van het raster. De doelcel moet leeg zijn, d.w.z. deze mag alleen het symbool '◌' bevatten.} \\
\texttt{* Als u probeert een object te verplaatsen naar een plek die niet leeg is, of als u probeert het buiten het raster te verplaatsen, moet ik u straffen. U krijgt een nieuwe poging.} \\
\texttt{* Voordat u een zet doet, moet u controleren of de doelplek leeg is, en geen letter, kader of lijn bevat!} \\
\\ 
\\ 
\texttt{**Einde van het spel**} \\
\\ 
\texttt{Als u denkt dat u het doel bereikt hebt om alle objecten uit te lijnen, kunt u de andere speler vragen het spel te beëindigen door `ZEG: klaar?` te sturen. Als de andere speler u vraagt om het spel te beëindigen en u antwoordt met `ZEG: klaar!`, zal het spel eindigen.} \\
\\ 
\texttt{Beide spelers winnen als het spel eindigt binnen 12 rondes, waarbij één ronde gedefinieerd wordt als twee spelers die elk een geldig commando sturen.} \\
\\ 
\texttt{**Score:**} \\
\\ 
\texttt{Hoe dichter de identieke objecten bij elkaar liggen op beide speelborden, hoe meer punten u krijgt. Strafpunten verminderen uw punten. Kunt u het record verslaan?} \\
\\ 
\\ 
\texttt{Stuur een bericht naar de andere speler om het spel te starten!} \\
            }
        }
    }
    & \\ \\

    \theutterance \stepcounter{utterance}  
    & \multicolumn{4}{p{0.6\linewidth}}{
        \cellcolor[rgb]{0.9,0.9,0.9}{
            \makecell[{{p{\linewidth}}}]{
                \texttt{\tiny{[P1$\rangle$GM]}}
                \texttt{ZEG: Hallo partner! Ik zie de objecten C, L en P op mijn bord. Laten we een strategie afspreken om ze te synchroniseren. Ik begin met het verplaatsen van objecten naar een bekende positie, bijvoorbeeld de linkerbovenhoek, en dan communiceren we onze posities. Ik verplaats nu C naar (1,1).} \\
            }
        }
    }
    & & \\ \\

    \theutterance \stepcounter{utterance}  
    & & & \multicolumn{2}{p{0.3\linewidth}}{
        \cellcolor[rgb]{0.9,0.9,0.9}{
            \makecell[{{p{\linewidth}}}]{
                \texttt{\tiny{[GM$|$GM]}}
                \texttt{ZEG: Hallo partner! Ik zie de objecten C, L en P op mijn bord. Laten we een strategie afspreken om ze te synchroniseren. Ik begin met het verplaatsen van objecten naar een bekende positie, bijvoorbeeld de linkerbovenhoek, en dan communiceren we onze posities. Ik verplaats nu C naar (1,1).} \\
            }
        }
    }
    & & \\ \\

    \theutterance \stepcounter{utterance}  
    & & \multicolumn{4}{p{0.6\linewidth}}{
        \cellcolor[rgb]{0.9,0.9,0.9}{
            \makecell[{{p{\linewidth}}}]{
                \texttt{\tiny{[GM$\rangle$P2]}}
                \texttt{Ik ben uw spelleider, en u speelt een samenwerkingsspel met het volgende raster als speelbord:} \\
\\ 
\texttt{```} \\
\texttt{ 1234567} \\
\texttt{╔══╤════╗} \\
\texttt{║◌◌│L◌◌C║ 1} \\
\texttt{╟──┤◌◌┌─╢ 2} \\
\texttt{║◌◌├──┤◌║ 3} \\
\texttt{╟──┘◌◌├─╢ 4} \\
\texttt{║◌◌P◌◌│◌║ 5} \\
\texttt{╟─┐◌┌─┼─╢ 6} \\
\texttt{║◌│◌│◌│◌║ 7} \\
\texttt{╚═╧═╧═╧═╝} \\
\\ 
\texttt{```} \\
\\ 
\texttt{* De bovenrand toont x{-}coördinaten die naar rechts oplopen, en de rechterrand y{-}coördinaten die naar beneden oplopen.} \\
\texttt{* De volgende objecten zijn willekeurig op uw raster geplaatst: 'C', 'L', 'P'.} \\
\\ 
\\ 
\texttt{De andere speler ziet een variatie van het speelbord, waarbij de objecten op verschillende willekeurige locaties zijn geplaatst. U kunt het bord van de andere speler niet zien en zij kunnen het uwe niet zien.} \\
\\ 
\\ 
\texttt{**Doel:**} \\
\\ 
\texttt{Beide spelers moeten de objecten op hun respectievelijke achtergrond zo verplaatsen dat identieke objecten op dezelfde coördinaten terechtkomen. U moet communiceren met de andere speler om een gemeenschappelijke doelconfiguratie overeen te komen.} \\
\\ 
\\ 
\texttt{**Regels:**} \\
\\ 
\texttt{* In elke beurt kunt u precies één van de volgende twee commando's sturen:} \\
\texttt{ 1. `ZEG: <BERICHT>`: om een bericht (alles tot de volgende regeleinde) naar de andere speler te sturen. Ik zal het doorsturen naar uw partner.} \\
\texttt{2. `BEWEEG: <OBJECT>, (<X>, <Y>)`: om een object naar een nieuwe positie te verplaatsen, waarbij `<X>` de kolom is en `<Y>` de rij. Ik zal u laten weten of uw zet succesvol was of niet.} \\
\texttt{* Als u zich niet aan het formaat houdt, of meerdere commando's tegelijk verstuurt, moet ik u straffen.} \\
\texttt{* Als beide spelers meer dan 8 strafpunten verzamelen, verliezen jullie allebei het spel.} \\
\texttt{* Het is van vitaal belang dat u met de andere speler communiceert over uw doelstatus! De *enige* manier waarop u uw strategie aan de andere speler kunt overbrengen is door gebruik te maken van het `ZEG: <BERICHT>` commando!} \\
\\ 
\\ 
\texttt{**Objecten verplaatsen**} \\
\\ 
\texttt{* U kunt alleen objecten verplaatsen naar cellen binnen de grenzen van het raster. De doelcel moet leeg zijn, d.w.z. deze mag alleen het symbool '◌' bevatten.} \\
\texttt{* Als u probeert een object te verplaatsen naar een plek die niet leeg is, of als u probeert het buiten het raster te verplaatsen, moet ik u straffen. U krijgt een nieuwe poging.} \\
\texttt{* Voordat u een zet doet, moet u controleren of de doelplek leeg is, en geen letter, kader of lijn bevat!} \\
\\ 
\\ 
\texttt{**Einde van het spel**} \\
\\ 
\texttt{Als u denkt dat u het doel bereikt hebt om alle objecten uit te lijnen, kunt u de andere speler vragen het spel te beëindigen door `ZEG: klaar?` te sturen. Als de andere speler u vraagt om het spel te beëindigen en u antwoordt met `ZEG: klaar!`, zal het spel eindigen.} \\
\\ 
\texttt{Beide spelers winnen als het spel eindigt binnen 12 rondes, waarbij één ronde gedefinieerd wordt als twee spelers die elk een geldig commando sturen.} \\
\\ 
\texttt{**Score:**} \\
\\ 
\texttt{Hoe dichter de identieke objecten bij elkaar liggen op beide speelborden, hoe meer punten u krijgt. Strafpunten verminderen uw punten. Kunt u het record verslaan?} \\
\\ 
\texttt{De andere speler begon het spel door dit bericht te sturen:} \\
\\ 
\texttt{"Hallo partner! Ik zie de objecten C, L en P op mijn bord. Laten we een strategie afspreken om ze te synchroniseren. Ik begin met het verplaatsen van objecten naar een bekende positie, bijvoorbeeld de linkerbovenhoek, en dan communiceren we onze posities. Ik verplaats nu C naar (1,1)."} \\
\\ 
\texttt{Wat is uw eerste opdracht?} \\
            }
        }
    }
    & \\ \\

    \theutterance \stepcounter{utterance}  
    & & & \multicolumn{4}{p{0.6\linewidth}}{
        \cellcolor[rgb]{0.9,0.9,0.9}{
            \makecell[{{p{\linewidth}}}]{
                \texttt{\tiny{[GM$\langle$P2]}}
                \texttt{Om het spel te winnen, moeten we eerst een gemeenschappelijke referentiekader creëren. Omdat we elkaars borden niet kunnen zien, is de enige manier om te synchroniseren dat we beide objecten naar dezelfde, bekende coördinaten verplaatsen (bijvoorbeeld de linkerbovenhoek (1,1)) en dan communiceren welke objecten waar staan.} \\
\\ 
\texttt{De andere speler heeft aangegeven C naar (1,1) te verplaatsen. Laten we hetzelfde doen voor ons bord, maar dan voor een ander object om verwarring te voorkomen, of we kunnen afspreken dat we allemaal naar (1,1) gaan en dan de volgorde vastleggen.} \\
\\ 
\texttt{Echter, een betere strategie is:} \\
\texttt{1. Verplaats alle objecten naar een vaste positie, bijvoorbeeld (1,1), (1,2), (1,3) in een vaste volgorde.} \\
\texttt{2. Communiceer welke objecten waar staan.} \\
\\ 
\texttt{Laten we beginnen met het verplaatsen van object 'C' naar (1,1) op ons bord, zodat we dezelfde positie delen als de partner (als zij C ook naar (1,1) verplaatst). Maar wacht, als zij C naar (1,1) verplaatst, en ik verplaats ook C naar (1,1), dan staan ze op dezelfde plek op *mijn* bord, maar niet noodzakelijk op *haar* bord.} \\
\\ 
\texttt{De sleutel is: **We moeten elkaars posities weten.**} \\
\\ 
\texttt{Laten we een eenvoudige strategie kiezen:} \\
\texttt{{-} Verplaats alle objecten naar de eerste drie cellen in de eerste rij: (1,1), (1,2), (1,3).} \\
\texttt{{-} Communiceer welke objecten waar staan.} \\
\\ 
\texttt{Op mijn bord staan:} \\
\texttt{{-} C op (7,1)} \\
\texttt{{-} L op (4,1)} \\
\texttt{{-} P op (3,5)} \\
\\ 
\texttt{Ik zal eerst C naar (1,1) verplaatsen.} \\
\\ 
\texttt{BEWEEG: C, (1,1)} \\
            }
        }
    }
     \\ \\

    \theutterance \stepcounter{utterance}  
    & & & \multicolumn{2}{p{0.3\linewidth}}{
        \cellcolor[rgb]{0.9,0.9,0.9}{
            \makecell[{{p{\linewidth}}}]{
                \texttt{\tiny{[GM$|$GM]}}
                \texttt{Om het spel te winnen, moeten we eerst een gemeenschappelijke referentiekader creëren. Omdat we elkaars borden niet kunnen zien, is de enige manier om te synchroniseren dat we beide objecten naar dezelfde, bekende coördinaten verplaatsen (bijvoorbeeld de linkerbovenhoek (1,1)) en dan communiceren welke objecten waar staan.} \\
\\ 
\texttt{De andere speler heeft aangegeven C naar (1,1) te verplaatsen. Laten we hetzelfde doen voor ons bord, maar dan voor een ander object om verwarring te voorkomen, of we kunnen afspreken dat we allemaal naar (1,1) gaan en dan de volgorde vastleggen.} \\
\\ 
\texttt{Echter, een betere strategie is:} \\
\texttt{1. Verplaats alle objecten naar een vaste positie, bijvoorbeeld (1,1), (1,2), (1,3) in een vaste volgorde.} \\
\texttt{2. Communiceer welke objecten waar staan.} \\
\\ 
\texttt{Laten we beginnen met het verplaatsen van object 'C' naar (1,1) op ons bord, zodat we dezelfde positie delen als de partner (als zij C ook naar (1,1) verplaatst). Maar wacht, als zij C naar (1,1) verplaatst, en ik verplaats ook C naar (1,1), dan staan ze op dezelfde plek op *mijn* bord, maar niet noodzakelijk op *haar* bord.} \\
\\ 
\texttt{De sleutel is: **We moeten elkaars posities weten.**} \\
\\ 
\texttt{Laten we een eenvoudige strategie kiezen:} \\
\texttt{{-} Verplaats alle objecten naar de eerste drie cellen in de eerste rij: (1,1), (1,2), (1,3).} \\
\texttt{{-} Communiceer welke objecten waar staan.} \\
\\ 
\texttt{Op mijn bord staan:} \\
\texttt{{-} C op (7,1)} \\
\texttt{{-} L op (4,1)} \\
\texttt{{-} P op (3,5)} \\
\\ 
\texttt{Ik zal eerst C naar (1,1) verplaatsen.} \\
\\ 
\texttt{BEWEEG: C, (1,1)} \\
            }
        }
    }
    & & \\ \\

    \theutterance \stepcounter{utterance}  
    & & & \multicolumn{2}{p{0.3\linewidth}}{
        \cellcolor[rgb]{0.9,0.9,0.9}{
            \makecell[{{p{\linewidth}}}]{
                \texttt{\tiny{[GM$|$GM]}}
                \texttt{Invalid format: head is not empty: 'Om het spel te winnen, moeten we eerst een gemeenschappelijke referentiekader creëren. Omdat we elkaars borden niet kunnen zien, is de enige manier om te synchroniseren dat we beide objecten naar dezelfde, bekende coördinaten verplaatsen (bijvoorbeeld de linkerbovenhoek (1,1)) en dan communiceren welke objecten waar staan.} \\
\\ 
\texttt{De andere speler heeft aangegeven C naar (1,1) te verplaatsen. Laten we hetzelfde doen voor ons bord, maar dan voor een ander object om verwarring te voorkomen, of we kunnen afspreken dat we allemaal naar (1,1) gaan en dan de volgorde vastleggen.} \\
\\ 
\texttt{Echter, een betere strategie is:} \\
\texttt{1. Verplaats alle objecten naar een vaste positie, bijvoorbeeld (1,1), (1,2), (1,3) in een vaste volgorde.} \\
\texttt{2. Communiceer welke objecten waar staan.} \\
\\ 
\texttt{Laten we beginnen met het verplaatsen van object 'C' naar (1,1) op ons bord, zodat we dezelfde positie delen als de partner (als zij C ook naar (1,1) verplaatst). Maar wacht, als zij C naar (1,1) verplaatst, en ik verplaats ook C naar (1,1), dan staan ze op dezelfde plek op *mijn* bord, maar niet noodzakelijk op *haar* bord.} \\
\\ 
\texttt{De sleutel is: **We moeten elkaars posities weten.**} \\
\\ 
\texttt{Laten we een eenvoudige strategie kiezen:} \\
\texttt{{-} Verplaats alle objecten naar de eerste drie cellen in de eerste rij: (1,1), (1,2), (1,3).} \\
\texttt{{-} Communiceer welke objecten waar staan.} \\
\\ 
\texttt{Op mijn bord staan:} \\
\texttt{{-} C op (7,1)} \\
\texttt{{-} L op (4,1)} \\
\texttt{{-} P op (3,5)} \\
\\ 
\texttt{Ik zal eerst C naar (1,1) verplaatsen.} \\
\\ 
\texttt{'} \\
            }
        }
    }
    & & \\ \\

    \theutterance \stepcounter{utterance}  
    & & \multicolumn{4}{p{0.6\linewidth}}{
        \cellcolor[rgb]{0.9,0.9,0.9}{
            \makecell[{{p{\linewidth}}}]{
                \texttt{\tiny{[GM$\rangle$P2]}}
                \texttt{Straf: Je bericht mag niets vóór het commando bevatten!} \\
\\ 
\texttt{Zorg ervoor dat je antwoord alleen `ZEG: <BERICHT>` of `BEWEEG: <OBJECT>, (<X>, <Y>)` bevat, en niets anders!} \\
\\ 
\texttt{Jullie hebben gezamenlijk 1 van 8 straffen verzameld.} \\
\texttt{Probeer het opnieuw!} \\
            }
        }
    }
    & \\ \\

    \theutterance \stepcounter{utterance}  
    & & & \multicolumn{4}{p{0.6\linewidth}}{
        \cellcolor[rgb]{0.9,0.9,0.9}{
            \makecell[{{p{\linewidth}}}]{
                \texttt{\tiny{[GM$\langle$P2]}}
                \texttt{BEWEEG: C, (1,1)} \\
            }
        }
    }
     \\ \\

    \theutterance \stepcounter{utterance}  
    & & & \multicolumn{2}{p{0.3\linewidth}}{
        \cellcolor[rgb]{0.9,0.9,0.9}{
            \makecell[{{p{\linewidth}}}]{
                \texttt{\tiny{[GM$|$GM]}}
                \texttt{BEWEEG: C, (1,1)} \\
            }
        }
    }
    & & \\ \\

    \theutterance \stepcounter{utterance}  
    & & & \multicolumn{2}{p{0.3\linewidth}}{
        \cellcolor[rgb]{0.9,0.9,0.9}{
            \makecell[{{p{\linewidth}}}]{
                \texttt{\tiny{[GM$|$GM]}}
                \texttt{C moved to (1, 1)} \\
            }
        }
    }
    & & \\ \\

    \theutterance \stepcounter{utterance}  
    & & \multicolumn{4}{p{0.6\linewidth}}{
        \cellcolor[rgb]{0.9,0.9,0.9}{
            \makecell[{{p{\linewidth}}}]{
                \texttt{\tiny{[P1$\langle$GM]}}
                \texttt{Je bericht is doorgestuurd naar de andere speler.} \\
\\ 
\\ 
\texttt{Je speelt momenteel ronde 1 van maximaal 12.} \\
\texttt{Jullie hebben gezamenlijk 1 van 8 straffen verzameld.} \\
\texttt{De andere speler heeft een voorwerp op zijn rooster verplaatst.} \\
\\ 
\\ 
\texttt{Wat is je volgende opdracht?} \\
            }
        }
    }
    & \\ \\

    \theutterance \stepcounter{utterance}  
    & \multicolumn{4}{p{0.6\linewidth}}{
        \cellcolor[rgb]{0.9,0.9,0.9}{
            \makecell[{{p{\linewidth}}}]{
                \texttt{\tiny{[P1$\rangle$GM]}}
                \texttt{ZEG: Ik heb C naar (1,1) verplaatst. Wat is jouw positie van C? En waar staan L en P bij jou? Laten we eerst C synchroniseren.} \\
            }
        }
    }
    & & \\ \\

    \theutterance \stepcounter{utterance}  
    & & & \multicolumn{2}{p{0.3\linewidth}}{
        \cellcolor[rgb]{0.9,0.9,0.9}{
            \makecell[{{p{\linewidth}}}]{
                \texttt{\tiny{[GM$|$GM]}}
                \texttt{ZEG: Ik heb C naar (1,1) verplaatst. Wat is jouw positie van C? En waar staan L en P bij jou? Laten we eerst C synchroniseren.} \\
            }
        }
    }
    & & \\ \\

    \theutterance \stepcounter{utterance}  
    & & \multicolumn{4}{p{0.6\linewidth}}{
        \cellcolor[rgb]{0.9,0.9,0.9}{
            \makecell[{{p{\linewidth}}}]{
                \texttt{\tiny{[GM$\rangle$P2]}}
                \texttt{Verplaatst 'C' naar (1,1) met succes. Uw bijgewerkte raster ziet er als volgt uit:} \\
\texttt{```} \\
\texttt{ 1234567} \\
\texttt{╔══╤════╗} \\
\texttt{║C◌│L◌◌◌║ 1} \\
\texttt{╟──┤◌◌┌─╢ 2} \\
\texttt{║◌◌├──┤◌║ 3} \\
\texttt{╟──┘◌◌├─╢ 4} \\
\texttt{║◌◌P◌◌│◌║ 5} \\
\texttt{╟─┐◌┌─┼─╢ 6} \\
\texttt{║◌│◌│◌│◌║ 7} \\
\texttt{╚═╧═╧═╧═╝} \\
\texttt{```} \\
\texttt{Je speelt momenteel ronde 2 van maximaal 12.} \\
\texttt{Jullie hebben gezamenlijk 1 van 8 straffen verzameld.} \\
\texttt{De andere speler stuurde dit bericht:} \\
\\ 
\\ 
\texttt{"Ik heb C naar (1,1) verplaatst. Wat is jouw positie van C? En waar staan L en P bij jou? Laten we eerst C synchroniseren."} \\
\\ 
\\ 
\texttt{Wat is je volgende opdracht?} \\
            }
        }
    }
    & \\ \\

    \theutterance \stepcounter{utterance}  
    & & & \multicolumn{4}{p{0.6\linewidth}}{
        \cellcolor[rgb]{0.9,0.9,0.9}{
            \makecell[{{p{\linewidth}}}]{
                \texttt{\tiny{[GM$\langle$P2]}}
                \texttt{ZEG: Ik heb C naar (1,1) verplaatst. L staat bij mij op (4,1) en P op (3,5). Laten we L ook naar (1,2) verplaatsen.} \\
            }
        }
    }
     \\ \\

    \theutterance \stepcounter{utterance}  
    & & & \multicolumn{2}{p{0.3\linewidth}}{
        \cellcolor[rgb]{0.9,0.9,0.9}{
            \makecell[{{p{\linewidth}}}]{
                \texttt{\tiny{[GM$|$GM]}}
                \texttt{ZEG: Ik heb C naar (1,1) verplaatst. L staat bij mij op (4,1) en P op (3,5). Laten we L ook naar (1,2) verplaatsen.} \\
            }
        }
    }
    & & \\ \\

    \theutterance \stepcounter{utterance}  
    & & \multicolumn{4}{p{0.6\linewidth}}{
        \cellcolor[rgb]{0.9,0.9,0.9}{
            \makecell[{{p{\linewidth}}}]{
                \texttt{\tiny{[P1$\langle$GM]}}
                \texttt{Je bericht is doorgestuurd naar de andere speler.} \\
\\ 
\\ 
\texttt{Je speelt momenteel ronde 2 van maximaal 12.} \\
\texttt{Jullie hebben gezamenlijk 1 van 8 straffen verzameld.} \\
\texttt{De andere speler stuurde dit bericht:} \\
\\ 
\\ 
\texttt{"Ik heb C naar (1,1) verplaatst. L staat bij mij op (4,1) en P op (3,5). Laten we L ook naar (1,2) verplaatsen."} \\
\\ 
\\ 
\texttt{Wat is je volgende opdracht?} \\
            }
        }
    }
    & \\ \\

    \theutterance \stepcounter{utterance}  
    & \multicolumn{4}{p{0.6\linewidth}}{
        \cellcolor[rgb]{0.9,0.9,0.9}{
            \makecell[{{p{\linewidth}}}]{
                \texttt{\tiny{[P1$\rangle$GM]}}
                \texttt{ZEG: Oké, ik zie dat jij C naar (1,1) hebt verplaatst. Ik heb C ook naar (1,1) verplaatst. Dus C is gesynchroniseerd op (1,1).} \\
\texttt{Jij zegt dat L bij jou op (4,1) staat en P op (3,5).} \\
\texttt{Ik zie L op mijn bord op (2,2) en P op (5,5).} \\
\texttt{Laten we L synchroniseren. Ik stel voor om L naar (1,2) te verplaatsen, zoals jij suggereerde. Ik verplaats L naar (1,2).} \\
            }
        }
    }
    & & \\ \\

    \theutterance \stepcounter{utterance}  
    & & & \multicolumn{2}{p{0.3\linewidth}}{
        \cellcolor[rgb]{0.9,0.9,0.9}{
            \makecell[{{p{\linewidth}}}]{
                \texttt{\tiny{[GM$|$GM]}}
                \texttt{ZEG: Oké, ik zie dat jij C naar (1,1) hebt verplaatst. Ik heb C ook naar (1,1) verplaatst. Dus C is gesynchroniseerd op (1,1).} \\
\texttt{Jij zegt dat L bij jou op (4,1) staat en P op (3,5).} \\
\texttt{Ik zie L op mijn bord op (2,2) en P op (5,5).} \\
\texttt{Laten we L synchroniseren. Ik stel voor om L naar (1,2) te verplaatsen, zoals jij suggereerde. Ik verplaats L naar (1,2).} \\
            }
        }
    }
    & & \\ \\

    \theutterance \stepcounter{utterance}  
    & & & \multicolumn{2}{p{0.3\linewidth}}{
        \cellcolor[rgb]{0.9,0.9,0.9}{
            \makecell[{{p{\linewidth}}}]{
                \texttt{\tiny{[GM$|$GM]}}
                \texttt{Invalid format: tail is not empty: '} \\
\texttt{Jij zegt dat L bij jou op (4,1) staat en P op (3,5).} \\
\texttt{Ik zie L op mijn bord op (2,2) en P op (5,5).} \\
\texttt{Laten we L synchroniseren. Ik stel voor om L naar (1,2) te verplaatsen, zoals jij suggereerde. Ik verplaats L naar (1,2).} \\
            }
        }
    }
    & & \\ \\

    \theutterance \stepcounter{utterance}  
    & & \multicolumn{4}{p{0.6\linewidth}}{
        \cellcolor[rgb]{0.9,0.9,0.9}{
            \makecell[{{p{\linewidth}}}]{
                \texttt{\tiny{[P1$\langle$GM]}}
                \texttt{Straf: Je bericht mag niets na het commando bevatten!} \\
\\ 
\texttt{Zorg ervoor dat je antwoord alleen `ZEG: <BERICHT>` of `BEWEEG: <OBJECT>, (<X>, <Y>)` bevat, en niets anders!} \\
\\ 
\texttt{Jullie hebben gezamenlijk 2 van 8 straffen verzameld.} \\
\texttt{Probeer het opnieuw!} \\
            }
        }
    }
    & \\ \\

    \theutterance \stepcounter{utterance}  
    & \multicolumn{4}{p{0.6\linewidth}}{
        \cellcolor[rgb]{0.9,0.9,0.9}{
            \makecell[{{p{\linewidth}}}]{
                \texttt{\tiny{[P1$\rangle$GM]}}
                \texttt{BEWEEG: L, (1,2)} \\
            }
        }
    }
    & & \\ \\

    \theutterance \stepcounter{utterance}  
    & & & \multicolumn{2}{p{0.3\linewidth}}{
        \cellcolor[rgb]{0.9,0.9,0.9}{
            \makecell[{{p{\linewidth}}}]{
                \texttt{\tiny{[GM$|$GM]}}
                \texttt{BEWEEG: L, (1,2)} \\
            }
        }
    }
    & & \\ \\

    \theutterance \stepcounter{utterance}  
    & & & \multicolumn{2}{p{0.3\linewidth}}{
        \cellcolor[rgb]{0.9,0.9,0.9}{
            \makecell[{{p{\linewidth}}}]{
                \texttt{\tiny{[GM$|$GM]}}
                \texttt{Straf: (1,2) is niet leeg, maar bevat '─'.} \\
\\ 
\texttt{Jullie hebben gezamenlijk 3 van 8 straffen verzameld.} \\
\texttt{Probeer het opnieuw!} \\
            }
        }
    }
    & & \\ \\

    \theutterance \stepcounter{utterance}  
    & & \multicolumn{4}{p{0.6\linewidth}}{
        \cellcolor[rgb]{0.9,0.9,0.9}{
            \makecell[{{p{\linewidth}}}]{
                \texttt{\tiny{[P1$\langle$GM]}}
                \texttt{Straf: (1,2) is niet leeg, maar bevat '─'.} \\
\\ 
\texttt{Jullie hebben gezamenlijk 3 van 8 straffen verzameld.} \\
\texttt{Probeer het opnieuw!} \\
            }
        }
    }
    & \\ \\

    \theutterance \stepcounter{utterance}  
    & \multicolumn{4}{p{0.6\linewidth}}{
        \cellcolor[rgb]{0.9,0.9,0.9}{
            \makecell[{{p{\linewidth}}}]{
                \texttt{\tiny{[P1$\rangle$GM]}}
                \texttt{ZEG: (1,2) is bezet door een lijn. Ik moet een andere plek zoeken voor L. Laten we eerst L en P synchroniseren op een lege plek. Ik zie dat (1,3) leeg is. Verplaats jij L naar (1,3)?} \\
            }
        }
    }
    & & \\ \\

    \theutterance \stepcounter{utterance}  
    & & & \multicolumn{2}{p{0.3\linewidth}}{
        \cellcolor[rgb]{0.9,0.9,0.9}{
            \makecell[{{p{\linewidth}}}]{
                \texttt{\tiny{[GM$|$GM]}}
                \texttt{ZEG: (1,2) is bezet door een lijn. Ik moet een andere plek zoeken voor L. Laten we eerst L en P synchroniseren op een lege plek. Ik zie dat (1,3) leeg is. Verplaats jij L naar (1,3)?} \\
            }
        }
    }
    & & \\ \\

    \theutterance \stepcounter{utterance}  
    & & \multicolumn{4}{p{0.6\linewidth}}{
        \cellcolor[rgb]{0.9,0.9,0.9}{
            \makecell[{{p{\linewidth}}}]{
                \texttt{\tiny{[GM$\rangle$P2]}}
                \texttt{Je bericht is doorgestuurd naar de andere speler.} \\
\\ 
\\ 
\texttt{Je speelt momenteel ronde 3 van maximaal 12.} \\
\texttt{Jullie hebben gezamenlijk 3 van 8 straffen verzameld.} \\
\texttt{De andere speler stuurde dit bericht:} \\
\\ 
\\ 
\texttt{"(1,2) is bezet door een lijn. Ik moet een andere plek zoeken voor L. Laten we eerst L en P synchroniseren op een lege plek. Ik zie dat (1,3) leeg is. Verplaats jij L naar (1,3)?"} \\
\\ 
\\ 
\texttt{Wat is je volgende opdracht?} \\
            }
        }
    }
    & \\ \\

    \theutterance \stepcounter{utterance}  
    & & & \multicolumn{4}{p{0.6\linewidth}}{
        \cellcolor[rgb]{0.9,0.9,0.9}{
            \makecell[{{p{\linewidth}}}]{
                \texttt{\tiny{[GM$\langle$P2]}}
                \texttt{ZEG: (1,3) is bij mij ook bezet door een lijn. Laten we L naar (1,2) verplaatsen, want bij mij is (1,2) leeg.} \\
            }
        }
    }
     \\ \\

    \theutterance \stepcounter{utterance}  
    & & & \multicolumn{2}{p{0.3\linewidth}}{
        \cellcolor[rgb]{0.9,0.9,0.9}{
            \makecell[{{p{\linewidth}}}]{
                \texttt{\tiny{[GM$|$GM]}}
                \texttt{ZEG: (1,3) is bij mij ook bezet door een lijn. Laten we L naar (1,2) verplaatsen, want bij mij is (1,2) leeg.} \\
            }
        }
    }
    & & \\ \\

    \theutterance \stepcounter{utterance}  
    & & \multicolumn{4}{p{0.6\linewidth}}{
        \cellcolor[rgb]{0.9,0.9,0.9}{
            \makecell[{{p{\linewidth}}}]{
                \texttt{\tiny{[P1$\langle$GM]}}
                \texttt{Je bericht is doorgestuurd naar de andere speler.} \\
\\ 
\\ 
\texttt{Je speelt momenteel ronde 3 van maximaal 12.} \\
\texttt{Jullie hebben gezamenlijk 3 van 8 straffen verzameld.} \\
\texttt{De andere speler stuurde dit bericht:} \\
\\ 
\\ 
\texttt{"(1,3) is bij mij ook bezet door een lijn. Laten we L naar (1,2) verplaatsen, want bij mij is (1,2) leeg."} \\
\\ 
\\ 
\texttt{Wat is je volgende opdracht?} \\
            }
        }
    }
    & \\ \\

    \theutterance \stepcounter{utterance}  
    & \multicolumn{4}{p{0.6\linewidth}}{
        \cellcolor[rgb]{0.9,0.9,0.9}{
            \makecell[{{p{\linewidth}}}]{
                \texttt{\tiny{[P1$\rangle$GM]}}
                \texttt{ZEG: Bij mij is (1,2) bezet door een lijn. Laten we L naar (3,1) verplaatsen, dat is leeg bij mij. Is (3,1) bij jou ook leeg?} \\
            }
        }
    }
    & & \\ \\

    \theutterance \stepcounter{utterance}  
    & & & \multicolumn{2}{p{0.3\linewidth}}{
        \cellcolor[rgb]{0.9,0.9,0.9}{
            \makecell[{{p{\linewidth}}}]{
                \texttt{\tiny{[GM$|$GM]}}
                \texttt{ZEG: Bij mij is (1,2) bezet door een lijn. Laten we L naar (3,1) verplaatsen, dat is leeg bij mij. Is (3,1) bij jou ook leeg?} \\
            }
        }
    }
    & & \\ \\

    \theutterance \stepcounter{utterance}  
    & & \multicolumn{4}{p{0.6\linewidth}}{
        \cellcolor[rgb]{0.9,0.9,0.9}{
            \makecell[{{p{\linewidth}}}]{
                \texttt{\tiny{[GM$\rangle$P2]}}
                \texttt{Je bericht is doorgestuurd naar de andere speler.} \\
\\ 
\\ 
\texttt{Je speelt momenteel ronde 4 van maximaal 12.} \\
\texttt{Jullie hebben gezamenlijk 3 van 8 straffen verzameld.} \\
\texttt{De andere speler stuurde dit bericht:} \\
\\ 
\\ 
\texttt{"Bij mij is (1,2) bezet door een lijn. Laten we L naar (3,1) verplaatsen, dat is leeg bij mij. Is (3,1) bij jou ook leeg?"} \\
\\ 
\\ 
\texttt{Wat is je volgende opdracht?} \\
            }
        }
    }
    & \\ \\

    \theutterance \stepcounter{utterance}  
    & & & \multicolumn{4}{p{0.6\linewidth}}{
        \cellcolor[rgb]{0.9,0.9,0.9}{
            \makecell[{{p{\linewidth}}}]{
                \texttt{\tiny{[GM$\langle$P2]}}
                \texttt{ZEG: Ja, (3,1) is bij mij leeg. Verplaats L naar (3,1).} \\
            }
        }
    }
     \\ \\

    \theutterance \stepcounter{utterance}  
    & & & \multicolumn{2}{p{0.3\linewidth}}{
        \cellcolor[rgb]{0.9,0.9,0.9}{
            \makecell[{{p{\linewidth}}}]{
                \texttt{\tiny{[GM$|$GM]}}
                \texttt{ZEG: Ja, (3,1) is bij mij leeg. Verplaats L naar (3,1).} \\
            }
        }
    }
    & & \\ \\

    \theutterance \stepcounter{utterance}  
    & & \multicolumn{4}{p{0.6\linewidth}}{
        \cellcolor[rgb]{0.9,0.9,0.9}{
            \makecell[{{p{\linewidth}}}]{
                \texttt{\tiny{[P1$\langle$GM]}}
                \texttt{Je bericht is doorgestuurd naar de andere speler.} \\
\\ 
\\ 
\texttt{Je speelt momenteel ronde 4 van maximaal 12.} \\
\texttt{Jullie hebben gezamenlijk 3 van 8 straffen verzameld.} \\
\texttt{De andere speler stuurde dit bericht:} \\
\\ 
\\ 
\texttt{"Ja, (3,1) is bij mij leeg. Verplaats L naar (3,1)."} \\
\\ 
\\ 
\texttt{Wat is je volgende opdracht?} \\
            }
        }
    }
    & \\ \\

    \theutterance \stepcounter{utterance}  
    & \multicolumn{4}{p{0.6\linewidth}}{
        \cellcolor[rgb]{0.9,0.9,0.9}{
            \makecell[{{p{\linewidth}}}]{
                \texttt{\tiny{[P1$\rangle$GM]}}
                \texttt{BEWEEG: L, (3,1)} \\
            }
        }
    }
    & & \\ \\

    \theutterance \stepcounter{utterance}  
    & & & \multicolumn{2}{p{0.3\linewidth}}{
        \cellcolor[rgb]{0.9,0.9,0.9}{
            \makecell[{{p{\linewidth}}}]{
                \texttt{\tiny{[GM$|$GM]}}
                \texttt{BEWEEG: L, (3,1)} \\
            }
        }
    }
    & & \\ \\

    \theutterance \stepcounter{utterance}  
    & & & \multicolumn{2}{p{0.3\linewidth}}{
        \cellcolor[rgb]{0.9,0.9,0.9}{
            \makecell[{{p{\linewidth}}}]{
                \texttt{\tiny{[GM$|$GM]}}
                \texttt{Straf: (3,1) is niet leeg, maar bevat '│'.} \\
\\ 
\texttt{Jullie hebben gezamenlijk 4 van 8 straffen verzameld.} \\
\texttt{Probeer het opnieuw!} \\
            }
        }
    }
    & & \\ \\

    \theutterance \stepcounter{utterance}  
    & & \multicolumn{4}{p{0.6\linewidth}}{
        \cellcolor[rgb]{0.9,0.9,0.9}{
            \makecell[{{p{\linewidth}}}]{
                \texttt{\tiny{[P1$\langle$GM]}}
                \texttt{Straf: (3,1) is niet leeg, maar bevat '│'.} \\
\\ 
\texttt{Jullie hebben gezamenlijk 4 van 8 straffen verzameld.} \\
\texttt{Probeer het opnieuw!} \\
            }
        }
    }
    & \\ \\

    \theutterance \stepcounter{utterance}  
    & \multicolumn{4}{p{0.6\linewidth}}{
        \cellcolor[rgb]{0.9,0.9,0.9}{
            \makecell[{{p{\linewidth}}}]{
                \texttt{\tiny{[P1$\rangle$GM]}}
                \texttt{ZEG: (3,1) is bezet door een lijn. Laten we L naar (1,3) verplaatsen. Is (1,3) bij jou leeg?} \\
            }
        }
    }
    & & \\ \\

    \theutterance \stepcounter{utterance}  
    & & & \multicolumn{2}{p{0.3\linewidth}}{
        \cellcolor[rgb]{0.9,0.9,0.9}{
            \makecell[{{p{\linewidth}}}]{
                \texttt{\tiny{[GM$|$GM]}}
                \texttt{ZEG: (3,1) is bezet door een lijn. Laten we L naar (1,3) verplaatsen. Is (1,3) bij jou leeg?} \\
            }
        }
    }
    & & \\ \\

    \theutterance \stepcounter{utterance}  
    & & \multicolumn{4}{p{0.6\linewidth}}{
        \cellcolor[rgb]{0.9,0.9,0.9}{
            \makecell[{{p{\linewidth}}}]{
                \texttt{\tiny{[GM$\rangle$P2]}}
                \texttt{Je bericht is doorgestuurd naar de andere speler.} \\
\\ 
\\ 
\texttt{Je speelt momenteel ronde 5 van maximaal 12.} \\
\texttt{Jullie hebben gezamenlijk 4 van 8 straffen verzameld.} \\
\texttt{De andere speler stuurde dit bericht:} \\
\\ 
\\ 
\texttt{"(3,1) is bezet door een lijn. Laten we L naar (1,3) verplaatsen. Is (1,3) bij jou leeg?"} \\
\\ 
\\ 
\texttt{Wat is je volgende opdracht?} \\
            }
        }
    }
    & \\ \\

    \theutterance \stepcounter{utterance}  
    & & & \multicolumn{4}{p{0.6\linewidth}}{
        \cellcolor[rgb]{0.9,0.9,0.9}{
            \makecell[{{p{\linewidth}}}]{
                \texttt{\tiny{[GM$\langle$P2]}}
                \texttt{ZEG: Nee, (1,3) is bij mij bezet door een lijn. Laten we L naar (2,1) verplaatsen. Is (2,1) bij jou leeg?} \\
            }
        }
    }
     \\ \\

    \theutterance \stepcounter{utterance}  
    & & & \multicolumn{2}{p{0.3\linewidth}}{
        \cellcolor[rgb]{0.9,0.9,0.9}{
            \makecell[{{p{\linewidth}}}]{
                \texttt{\tiny{[GM$|$GM]}}
                \texttt{ZEG: Nee, (1,3) is bij mij bezet door een lijn. Laten we L naar (2,1) verplaatsen. Is (2,1) bij jou leeg?} \\
            }
        }
    }
    & & \\ \\

    \theutterance \stepcounter{utterance}  
    & & \multicolumn{4}{p{0.6\linewidth}}{
        \cellcolor[rgb]{0.9,0.9,0.9}{
            \makecell[{{p{\linewidth}}}]{
                \texttt{\tiny{[P1$\langle$GM]}}
                \texttt{Je bericht is doorgestuurd naar de andere speler.} \\
\\ 
\\ 
\texttt{Je speelt momenteel ronde 5 van maximaal 12.} \\
\texttt{Jullie hebben gezamenlijk 4 van 8 straffen verzameld.} \\
\texttt{De andere speler stuurde dit bericht:} \\
\\ 
\\ 
\texttt{"Nee, (1,3) is bij mij bezet door een lijn. Laten we L naar (2,1) verplaatsen. Is (2,1) bij jou leeg?"} \\
\\ 
\\ 
\texttt{Wat is je volgende opdracht?} \\
            }
        }
    }
    & \\ \\

    \theutterance \stepcounter{utterance}  
    & \multicolumn{4}{p{0.6\linewidth}}{
        \cellcolor[rgb]{0.9,0.9,0.9}{
            \makecell[{{p{\linewidth}}}]{
                \texttt{\tiny{[P1$\rangle$GM]}}
                \texttt{ZEG: (2,1) is bij mij ook bezet. Laten we L naar (4,1) verplaatsen. Is (4,1) bij jou leeg?} \\
            }
        }
    }
    & & \\ \\

    \theutterance \stepcounter{utterance}  
    & & & \multicolumn{2}{p{0.3\linewidth}}{
        \cellcolor[rgb]{0.9,0.9,0.9}{
            \makecell[{{p{\linewidth}}}]{
                \texttt{\tiny{[GM$|$GM]}}
                \texttt{ZEG: (2,1) is bij mij ook bezet. Laten we L naar (4,1) verplaatsen. Is (4,1) bij jou leeg?} \\
            }
        }
    }
    & & \\ \\

    \theutterance \stepcounter{utterance}  
    & & \multicolumn{4}{p{0.6\linewidth}}{
        \cellcolor[rgb]{0.9,0.9,0.9}{
            \makecell[{{p{\linewidth}}}]{
                \texttt{\tiny{[GM$\rangle$P2]}}
                \texttt{Je bericht is doorgestuurd naar de andere speler.} \\
\\ 
\\ 
\texttt{Je speelt momenteel ronde 6 van maximaal 12.} \\
\texttt{Jullie hebben gezamenlijk 4 van 8 straffen verzameld.} \\
\texttt{De andere speler stuurde dit bericht:} \\
\\ 
\\ 
\texttt{"(2,1) is bij mij ook bezet. Laten we L naar (4,1) verplaatsen. Is (4,1) bij jou leeg?"} \\
\\ 
\\ 
\texttt{Wat is je volgende opdracht?} \\
            }
        }
    }
    & \\ \\

    \theutterance \stepcounter{utterance}  
    & & & \multicolumn{4}{p{0.6\linewidth}}{
        \cellcolor[rgb]{0.9,0.9,0.9}{
            \makecell[{{p{\linewidth}}}]{
                \texttt{\tiny{[GM$\langle$P2]}}
                \texttt{ZEG: Nee, (4,1) is bij mij bezet door L. Laten we L naar (5,1) verplaatsen. Is (5,1) bij jou leeg?} \\
            }
        }
    }
     \\ \\

    \theutterance \stepcounter{utterance}  
    & & & \multicolumn{2}{p{0.3\linewidth}}{
        \cellcolor[rgb]{0.9,0.9,0.9}{
            \makecell[{{p{\linewidth}}}]{
                \texttt{\tiny{[GM$|$GM]}}
                \texttt{ZEG: Nee, (4,1) is bij mij bezet door L. Laten we L naar (5,1) verplaatsen. Is (5,1) bij jou leeg?} \\
            }
        }
    }
    & & \\ \\

    \theutterance \stepcounter{utterance}  
    & & \multicolumn{4}{p{0.6\linewidth}}{
        \cellcolor[rgb]{0.9,0.9,0.9}{
            \makecell[{{p{\linewidth}}}]{
                \texttt{\tiny{[P1$\langle$GM]}}
                \texttt{Je bericht is doorgestuurd naar de andere speler.} \\
\\ 
\\ 
\texttt{Je speelt momenteel ronde 6 van maximaal 12.} \\
\texttt{Jullie hebben gezamenlijk 4 van 8 straffen verzameld.} \\
\texttt{De andere speler stuurde dit bericht:} \\
\\ 
\\ 
\texttt{"Nee, (4,1) is bij mij bezet door L. Laten we L naar (5,1) verplaatsen. Is (5,1) bij jou leeg?"} \\
\\ 
\\ 
\texttt{Wat is je volgende opdracht?} \\
            }
        }
    }
    & \\ \\

    \theutterance \stepcounter{utterance}  
    & \multicolumn{4}{p{0.6\linewidth}}{
        \cellcolor[rgb]{0.9,0.9,0.9}{
            \makecell[{{p{\linewidth}}}]{
                \texttt{\tiny{[P1$\rangle$GM]}}
                \texttt{ZEG: (5,1) is bij mij bezet door P. Laten we L naar (6,1) verplaatsen. Is (6,1) bij jou leeg?} \\
            }
        }
    }
    & & \\ \\

    \theutterance \stepcounter{utterance}  
    & & & \multicolumn{2}{p{0.3\linewidth}}{
        \cellcolor[rgb]{0.9,0.9,0.9}{
            \makecell[{{p{\linewidth}}}]{
                \texttt{\tiny{[GM$|$GM]}}
                \texttt{ZEG: (5,1) is bij mij bezet door P. Laten we L naar (6,1) verplaatsen. Is (6,1) bij jou leeg?} \\
            }
        }
    }
    & & \\ \\

    \theutterance \stepcounter{utterance}  
    & & \multicolumn{4}{p{0.6\linewidth}}{
        \cellcolor[rgb]{0.9,0.9,0.9}{
            \makecell[{{p{\linewidth}}}]{
                \texttt{\tiny{[GM$\rangle$P2]}}
                \texttt{Je bericht is doorgestuurd naar de andere speler.} \\
\\ 
\\ 
\texttt{Je speelt momenteel ronde 7 van maximaal 12.} \\
\texttt{Jullie hebben gezamenlijk 4 van 8 straffen verzameld.} \\
\texttt{De andere speler stuurde dit bericht:} \\
\\ 
\\ 
\texttt{"(5,1) is bij mij bezet door P. Laten we L naar (6,1) verplaatsen. Is (6,1) bij jou leeg?"} \\
\\ 
\\ 
\texttt{Wat is je volgende opdracht?} \\
            }
        }
    }
    & \\ \\

    \theutterance \stepcounter{utterance}  
    & & & \multicolumn{4}{p{0.6\linewidth}}{
        \cellcolor[rgb]{0.9,0.9,0.9}{
            \makecell[{{p{\linewidth}}}]{
                \texttt{\tiny{[GM$\langle$P2]}}
                \texttt{ZEG: Ja, (6,1) is bij mij leeg. Verplaats L naar (6,1).} \\
            }
        }
    }
     \\ \\

    \theutterance \stepcounter{utterance}  
    & & & \multicolumn{2}{p{0.3\linewidth}}{
        \cellcolor[rgb]{0.9,0.9,0.9}{
            \makecell[{{p{\linewidth}}}]{
                \texttt{\tiny{[GM$|$GM]}}
                \texttt{ZEG: Ja, (6,1) is bij mij leeg. Verplaats L naar (6,1).} \\
            }
        }
    }
    & & \\ \\

    \theutterance \stepcounter{utterance}  
    & & \multicolumn{4}{p{0.6\linewidth}}{
        \cellcolor[rgb]{0.9,0.9,0.9}{
            \makecell[{{p{\linewidth}}}]{
                \texttt{\tiny{[P1$\langle$GM]}}
                \texttt{Je bericht is doorgestuurd naar de andere speler.} \\
\\ 
\\ 
\texttt{Je speelt momenteel ronde 7 van maximaal 12.} \\
\texttt{Jullie hebben gezamenlijk 4 van 8 straffen verzameld.} \\
\texttt{De andere speler stuurde dit bericht:} \\
\\ 
\\ 
\texttt{"Ja, (6,1) is bij mij leeg. Verplaats L naar (6,1)."} \\
\\ 
\\ 
\texttt{Wat is je volgende opdracht?} \\
            }
        }
    }
    & \\ \\

    \theutterance \stepcounter{utterance}  
    & \multicolumn{4}{p{0.6\linewidth}}{
        \cellcolor[rgb]{0.9,0.9,0.9}{
            \makecell[{{p{\linewidth}}}]{
                \texttt{\tiny{[P1$\rangle$GM]}}
                \texttt{BEWEEG: L, (6,1)} \\
            }
        }
    }
    & & \\ \\

    \theutterance \stepcounter{utterance}  
    & & & \multicolumn{2}{p{0.3\linewidth}}{
        \cellcolor[rgb]{0.9,0.9,0.9}{
            \makecell[{{p{\linewidth}}}]{
                \texttt{\tiny{[GM$|$GM]}}
                \texttt{BEWEEG: L, (6,1)} \\
            }
        }
    }
    & & \\ \\

    \theutterance \stepcounter{utterance}  
    & & & \multicolumn{2}{p{0.3\linewidth}}{
        \cellcolor[rgb]{0.9,0.9,0.9}{
            \makecell[{{p{\linewidth}}}]{
                \texttt{\tiny{[GM$|$GM]}}
                \texttt{L moved to (6, 1)} \\
            }
        }
    }
    & & \\ \\

    \theutterance \stepcounter{utterance}  
    & & \multicolumn{4}{p{0.6\linewidth}}{
        \cellcolor[rgb]{0.9,0.9,0.9}{
            \makecell[{{p{\linewidth}}}]{
                \texttt{\tiny{[GM$\rangle$P2]}}
                \texttt{Je bericht is doorgestuurd naar de andere speler.} \\
\\ 
\\ 
\texttt{Je speelt momenteel ronde 8 van maximaal 12.} \\
\texttt{Jullie hebben gezamenlijk 4 van 8 straffen verzameld.} \\
\texttt{De andere speler heeft een voorwerp op zijn rooster verplaatst.} \\
\\ 
\\ 
\texttt{Wat is je volgende opdracht?} \\
            }
        }
    }
    & \\ \\

    \theutterance \stepcounter{utterance}  
    & & & \multicolumn{4}{p{0.6\linewidth}}{
        \cellcolor[rgb]{0.9,0.9,0.9}{
            \makecell[{{p{\linewidth}}}]{
                \texttt{\tiny{[GM$\langle$P2]}}
                \texttt{BEWEEG: L, (6,1)} \\
            }
        }
    }
     \\ \\

    \theutterance \stepcounter{utterance}  
    & & & \multicolumn{2}{p{0.3\linewidth}}{
        \cellcolor[rgb]{0.9,0.9,0.9}{
            \makecell[{{p{\linewidth}}}]{
                \texttt{\tiny{[GM$|$GM]}}
                \texttt{BEWEEG: L, (6,1)} \\
            }
        }
    }
    & & \\ \\

    \theutterance \stepcounter{utterance}  
    & & & \multicolumn{2}{p{0.3\linewidth}}{
        \cellcolor[rgb]{0.9,0.9,0.9}{
            \makecell[{{p{\linewidth}}}]{
                \texttt{\tiny{[GM$|$GM]}}
                \texttt{L moved to (6, 1)} \\
            }
        }
    }
    & & \\ \\

    \theutterance \stepcounter{utterance}  
    & & \multicolumn{4}{p{0.6\linewidth}}{
        \cellcolor[rgb]{0.9,0.9,0.9}{
            \makecell[{{p{\linewidth}}}]{
                \texttt{\tiny{[P1$\langle$GM]}}
                \texttt{Verplaatst 'L' naar (6,1) met succes. Uw bijgewerkte raster ziet er als volgt uit:} \\
\texttt{```} \\
\texttt{ 1234567} \\
\texttt{╔══╤════╗} \\
\texttt{║◌◌│C◌L◌║ 1} \\
\texttt{╟──┤◌◌┌─╢ 2} \\
\texttt{║◌◌├──┤◌║ 3} \\
\texttt{╟──┘◌◌├─╢ 4} \\
\texttt{║◌◌P◌◌│◌║ 5} \\
\texttt{╟─┐◌┌─┼─╢ 6} \\
\texttt{║◌│◌│◌│◌║ 7} \\
\texttt{╚═╧═╧═╧═╝} \\
\texttt{```} \\
\texttt{Je speelt momenteel ronde 8 van maximaal 12.} \\
\texttt{Jullie hebben gezamenlijk 4 van 8 straffen verzameld.} \\
\texttt{De andere speler heeft een voorwerp op zijn rooster verplaatst.} \\
\\ 
\\ 
\texttt{Wat is je volgende opdracht?} \\
            }
        }
    }
    & \\ \\

    \theutterance \stepcounter{utterance}  
    & \multicolumn{4}{p{0.6\linewidth}}{
        \cellcolor[rgb]{0.9,0.9,0.9}{
            \makecell[{{p{\linewidth}}}]{
                \texttt{\tiny{[P1$\rangle$GM]}}
                \texttt{ZEG: L is nu op (6,1) bij mij. Is L bij jou ook op (6,1)? Nu moeten we C en P synchroniseren. Waar staat C bij jou?} \\
            }
        }
    }
    & & \\ \\

    \theutterance \stepcounter{utterance}  
    & & & \multicolumn{2}{p{0.3\linewidth}}{
        \cellcolor[rgb]{0.9,0.9,0.9}{
            \makecell[{{p{\linewidth}}}]{
                \texttt{\tiny{[GM$|$GM]}}
                \texttt{ZEG: L is nu op (6,1) bij mij. Is L bij jou ook op (6,1)? Nu moeten we C en P synchroniseren. Waar staat C bij jou?} \\
            }
        }
    }
    & & \\ \\

    \theutterance \stepcounter{utterance}  
    & & \multicolumn{4}{p{0.6\linewidth}}{
        \cellcolor[rgb]{0.9,0.9,0.9}{
            \makecell[{{p{\linewidth}}}]{
                \texttt{\tiny{[GM$\rangle$P2]}}
                \texttt{Verplaatst 'L' naar (6,1) met succes. Uw bijgewerkte raster ziet er als volgt uit:} \\
\texttt{```} \\
\texttt{ 1234567} \\
\texttt{╔══╤════╗} \\
\texttt{║C◌│◌◌L◌║ 1} \\
\texttt{╟──┤◌◌┌─╢ 2} \\
\texttt{║◌◌├──┤◌║ 3} \\
\texttt{╟──┘◌◌├─╢ 4} \\
\texttt{║◌◌P◌◌│◌║ 5} \\
\texttt{╟─┐◌┌─┼─╢ 6} \\
\texttt{║◌│◌│◌│◌║ 7} \\
\texttt{╚═╧═╧═╧═╝} \\
\texttt{```} \\
\texttt{Je speelt momenteel ronde 9 van maximaal 12.} \\
\texttt{Jullie hebben gezamenlijk 4 van 8 straffen verzameld.} \\
\texttt{De andere speler stuurde dit bericht:} \\
\\ 
\\ 
\texttt{"L is nu op (6,1) bij mij. Is L bij jou ook op (6,1)? Nu moeten we C en P synchroniseren. Waar staat C bij jou?"} \\
\\ 
\\ 
\texttt{Wat is je volgende opdracht?} \\
            }
        }
    }
    & \\ \\

    \theutterance \stepcounter{utterance}  
    & & & \multicolumn{4}{p{0.6\linewidth}}{
        \cellcolor[rgb]{0.9,0.9,0.9}{
            \makecell[{{p{\linewidth}}}]{
                \texttt{\tiny{[GM$\langle$P2]}}
                \texttt{ZEG: Ja, L staat bij mij ook op (6,1). C staat bij mij op (1,1). Laten we P synchroniseren. Waar staat P bij jou?} \\
            }
        }
    }
     \\ \\

    \theutterance \stepcounter{utterance}  
    & & & \multicolumn{2}{p{0.3\linewidth}}{
        \cellcolor[rgb]{0.9,0.9,0.9}{
            \makecell[{{p{\linewidth}}}]{
                \texttt{\tiny{[GM$|$GM]}}
                \texttt{ZEG: Ja, L staat bij mij ook op (6,1). C staat bij mij op (1,1). Laten we P synchroniseren. Waar staat P bij jou?} \\
            }
        }
    }
    & & \\ \\

    \theutterance \stepcounter{utterance}  
    & & \multicolumn{4}{p{0.6\linewidth}}{
        \cellcolor[rgb]{0.9,0.9,0.9}{
            \makecell[{{p{\linewidth}}}]{
                \texttt{\tiny{[P1$\langle$GM]}}
                \texttt{Je bericht is doorgestuurd naar de andere speler.} \\
\\ 
\\ 
\texttt{Je speelt momenteel ronde 9 van maximaal 12.} \\
\texttt{Jullie hebben gezamenlijk 4 van 8 straffen verzameld.} \\
\texttt{De andere speler stuurde dit bericht:} \\
\\ 
\\ 
\texttt{"Ja, L staat bij mij ook op (6,1). C staat bij mij op (1,1). Laten we P synchroniseren. Waar staat P bij jou?"} \\
\\ 
\\ 
\texttt{Wat is je volgende opdracht?} \\
            }
        }
    }
    & \\ \\

    \theutterance \stepcounter{utterance}  
    & \multicolumn{4}{p{0.6\linewidth}}{
        \cellcolor[rgb]{0.9,0.9,0.9}{
            \makecell[{{p{\linewidth}}}]{
                \texttt{\tiny{[P1$\rangle$GM]}}
                \texttt{ZEG: P staat bij mij op (5,5). Laten we P naar (5,1) verplaatsen. Is (5,1) bij jou leeg?} \\
            }
        }
    }
    & & \\ \\

    \theutterance \stepcounter{utterance}  
    & & & \multicolumn{2}{p{0.3\linewidth}}{
        \cellcolor[rgb]{0.9,0.9,0.9}{
            \makecell[{{p{\linewidth}}}]{
                \texttt{\tiny{[GM$|$GM]}}
                \texttt{ZEG: P staat bij mij op (5,5). Laten we P naar (5,1) verplaatsen. Is (5,1) bij jou leeg?} \\
            }
        }
    }
    & & \\ \\

    \theutterance \stepcounter{utterance}  
    & & \multicolumn{4}{p{0.6\linewidth}}{
        \cellcolor[rgb]{0.9,0.9,0.9}{
            \makecell[{{p{\linewidth}}}]{
                \texttt{\tiny{[GM$\rangle$P2]}}
                \texttt{Je bericht is doorgestuurd naar de andere speler.} \\
\\ 
\\ 
\texttt{Je speelt momenteel ronde 10 van maximaal 12.} \\
\texttt{Jullie hebben gezamenlijk 4 van 8 straffen verzameld.} \\
\texttt{De andere speler stuurde dit bericht:} \\
\\ 
\\ 
\texttt{"P staat bij mij op (5,5). Laten we P naar (5,1) verplaatsen. Is (5,1) bij jou leeg?"} \\
\\ 
\\ 
\texttt{Wat is je volgende opdracht?} \\
            }
        }
    }
    & \\ \\

    \theutterance \stepcounter{utterance}  
    & & & \multicolumn{4}{p{0.6\linewidth}}{
        \cellcolor[rgb]{0.9,0.9,0.9}{
            \makecell[{{p{\linewidth}}}]{
                \texttt{\tiny{[GM$\langle$P2]}}
                \texttt{ZEG: Nee, (5,1) is bij mij bezet door P. Laten we P naar (5,2) verplaatsen. Is (5,2) bij jou leeg?} \\
            }
        }
    }
     \\ \\

    \theutterance \stepcounter{utterance}  
    & & & \multicolumn{2}{p{0.3\linewidth}}{
        \cellcolor[rgb]{0.9,0.9,0.9}{
            \makecell[{{p{\linewidth}}}]{
                \texttt{\tiny{[GM$|$GM]}}
                \texttt{ZEG: Nee, (5,1) is bij mij bezet door P. Laten we P naar (5,2) verplaatsen. Is (5,2) bij jou leeg?} \\
            }
        }
    }
    & & \\ \\

    \theutterance \stepcounter{utterance}  
    & & \multicolumn{4}{p{0.6\linewidth}}{
        \cellcolor[rgb]{0.9,0.9,0.9}{
            \makecell[{{p{\linewidth}}}]{
                \texttt{\tiny{[P1$\langle$GM]}}
                \texttt{Je bericht is doorgestuurd naar de andere speler.} \\
\\ 
\\ 
\texttt{Je speelt momenteel ronde 10 van maximaal 12.} \\
\texttt{Jullie hebben gezamenlijk 4 van 8 straffen verzameld.} \\
\texttt{De andere speler stuurde dit bericht:} \\
\\ 
\\ 
\texttt{"Nee, (5,1) is bij mij bezet door P. Laten we P naar (5,2) verplaatsen. Is (5,2) bij jou leeg?"} \\
\\ 
\\ 
\texttt{Wat is je volgende opdracht?} \\
            }
        }
    }
    & \\ \\

    \theutterance \stepcounter{utterance}  
    & \multicolumn{4}{p{0.6\linewidth}}{
        \cellcolor[rgb]{0.9,0.9,0.9}{
            \makecell[{{p{\linewidth}}}]{
                \texttt{\tiny{[P1$\rangle$GM]}}
                \texttt{ZEG: (5,2) is bij mij bezet door een lijn. Laten we P naar (5,3) verplaatsen. Is (5,3) bij jou leeg?} \\
            }
        }
    }
    & & \\ \\

    \theutterance \stepcounter{utterance}  
    & & & \multicolumn{2}{p{0.3\linewidth}}{
        \cellcolor[rgb]{0.9,0.9,0.9}{
            \makecell[{{p{\linewidth}}}]{
                \texttt{\tiny{[GM$|$GM]}}
                \texttt{ZEG: (5,2) is bij mij bezet door een lijn. Laten we P naar (5,3) verplaatsen. Is (5,3) bij jou leeg?} \\
            }
        }
    }
    & & \\ \\

    \theutterance \stepcounter{utterance}  
    & & \multicolumn{4}{p{0.6\linewidth}}{
        \cellcolor[rgb]{0.9,0.9,0.9}{
            \makecell[{{p{\linewidth}}}]{
                \texttt{\tiny{[GM$\rangle$P2]}}
                \texttt{Je bericht is doorgestuurd naar de andere speler.} \\
\\ 
\\ 
\texttt{Je speelt momenteel ronde 11 van maximaal 12.} \\
\texttt{Jullie hebben gezamenlijk 4 van 8 straffen verzameld.} \\
\texttt{De andere speler stuurde dit bericht:} \\
\\ 
\\ 
\texttt{"(5,2) is bij mij bezet door een lijn. Laten we P naar (5,3) verplaatsen. Is (5,3) bij jou leeg?"} \\
\\ 
\\ 
\texttt{Wat is je volgende opdracht?} \\
            }
        }
    }
    & \\ \\

    \theutterance \stepcounter{utterance}  
    & & & \multicolumn{4}{p{0.6\linewidth}}{
        \cellcolor[rgb]{0.9,0.9,0.9}{
            \makecell[{{p{\linewidth}}}]{
                \texttt{\tiny{[GM$\langle$P2]}}
                \texttt{ZEG: Ja, (5,3) is bij mij leeg. Verplaats P naar (5,3).} \\
            }
        }
    }
     \\ \\

    \theutterance \stepcounter{utterance}  
    & & & \multicolumn{2}{p{0.3\linewidth}}{
        \cellcolor[rgb]{0.9,0.9,0.9}{
            \makecell[{{p{\linewidth}}}]{
                \texttt{\tiny{[GM$|$GM]}}
                \texttt{ZEG: Ja, (5,3) is bij mij leeg. Verplaats P naar (5,3).} \\
            }
        }
    }
    & & \\ \\

    \theutterance \stepcounter{utterance}  
    & & \multicolumn{4}{p{0.6\linewidth}}{
        \cellcolor[rgb]{0.9,0.9,0.9}{
            \makecell[{{p{\linewidth}}}]{
                \texttt{\tiny{[P1$\langle$GM]}}
                \texttt{Je bericht is doorgestuurd naar de andere speler.} \\
\\ 
\\ 
\texttt{Je speelt momenteel ronde 11 van maximaal 12.} \\
\texttt{Jullie hebben gezamenlijk 4 van 8 straffen verzameld.} \\
\texttt{De andere speler stuurde dit bericht:} \\
\\ 
\\ 
\texttt{"Ja, (5,3) is bij mij leeg. Verplaats P naar (5,3)."} \\
\\ 
\\ 
\texttt{Wat is je volgende opdracht?} \\
            }
        }
    }
    & \\ \\

    \theutterance \stepcounter{utterance}  
    & \multicolumn{4}{p{0.6\linewidth}}{
        \cellcolor[rgb]{0.9,0.9,0.9}{
            \makecell[{{p{\linewidth}}}]{
                \texttt{\tiny{[P1$\rangle$GM]}}
                \texttt{BEWEEG: P, (5,3)} \\
            }
        }
    }
    & & \\ \\

    \theutterance \stepcounter{utterance}  
    & & & \multicolumn{2}{p{0.3\linewidth}}{
        \cellcolor[rgb]{0.9,0.9,0.9}{
            \makecell[{{p{\linewidth}}}]{
                \texttt{\tiny{[GM$|$GM]}}
                \texttt{BEWEEG: P, (5,3)} \\
            }
        }
    }
    & & \\ \\

    \theutterance \stepcounter{utterance}  
    & & & \multicolumn{2}{p{0.3\linewidth}}{
        \cellcolor[rgb]{0.9,0.9,0.9}{
            \makecell[{{p{\linewidth}}}]{
                \texttt{\tiny{[GM$|$GM]}}
                \texttt{Straf: (5,3) is niet leeg, maar bevat '─'.} \\
\\ 
\texttt{Jullie hebben gezamenlijk 5 van 8 straffen verzameld.} \\
\texttt{Probeer het opnieuw!} \\
            }
        }
    }
    & & \\ \\

    \theutterance \stepcounter{utterance}  
    & & \multicolumn{4}{p{0.6\linewidth}}{
        \cellcolor[rgb]{0.9,0.9,0.9}{
            \makecell[{{p{\linewidth}}}]{
                \texttt{\tiny{[P1$\langle$GM]}}
                \texttt{Straf: (5,3) is niet leeg, maar bevat '─'.} \\
\\ 
\texttt{Jullie hebben gezamenlijk 5 van 8 straffen verzameld.} \\
\texttt{Probeer het opnieuw!} \\
            }
        }
    }
    & \\ \\

    \theutterance \stepcounter{utterance}  
    & \multicolumn{4}{p{0.6\linewidth}}{
        \cellcolor[rgb]{0.9,0.9,0.9}{
            \makecell[{{p{\linewidth}}}]{
                \texttt{\tiny{[P1$\rangle$GM]}}
                \texttt{ZEG: (5,3) is bezet door een lijn. Laten we P naar (5,4) verplaatsen. Is (5,4) bij jou leeg?} \\
            }
        }
    }
    & & \\ \\

    \theutterance \stepcounter{utterance}  
    & & & \multicolumn{2}{p{0.3\linewidth}}{
        \cellcolor[rgb]{0.9,0.9,0.9}{
            \makecell[{{p{\linewidth}}}]{
                \texttt{\tiny{[GM$|$GM]}}
                \texttt{ZEG: (5,3) is bezet door een lijn. Laten we P naar (5,4) verplaatsen. Is (5,4) bij jou leeg?} \\
            }
        }
    }
    & & \\ \\

    \theutterance \stepcounter{utterance}  
    & & \multicolumn{4}{p{0.6\linewidth}}{
        \cellcolor[rgb]{0.9,0.9,0.9}{
            \makecell[{{p{\linewidth}}}]{
                \texttt{\tiny{[GM$\rangle$P2]}}
                \texttt{Je bericht is doorgestuurd naar de andere speler.} \\
\\ 
\\ 
\texttt{Je speelt momenteel ronde 12 van maximaal 12.} \\
\texttt{Jullie hebben gezamenlijk 5 van 8 straffen verzameld.} \\
\texttt{De andere speler stuurde dit bericht:} \\
\\ 
\\ 
\texttt{"(5,3) is bezet door een lijn. Laten we P naar (5,4) verplaatsen. Is (5,4) bij jou leeg?"} \\
\\ 
\\ 
\texttt{Wat is je volgende opdracht?} \\
            }
        }
    }
    & \\ \\

    \theutterance \stepcounter{utterance}  
    & & & \multicolumn{4}{p{0.6\linewidth}}{
        \cellcolor[rgb]{0.9,0.9,0.9}{
            \makecell[{{p{\linewidth}}}]{
                \texttt{\tiny{[GM$\langle$P2]}}
                \texttt{ZEG: Ja, (5,4) is bij mij leeg. Verplaats P naar (5,4).} \\
            }
        }
    }
     \\ \\

    \theutterance \stepcounter{utterance}  
    & & & \multicolumn{2}{p{0.3\linewidth}}{
        \cellcolor[rgb]{0.9,0.9,0.9}{
            \makecell[{{p{\linewidth}}}]{
                \texttt{\tiny{[GM$|$GM]}}
                \texttt{ZEG: Ja, (5,4) is bij mij leeg. Verplaats P naar (5,4).} \\
            }
        }
    }
    & & \\ \\

    \theutterance \stepcounter{utterance}  
    & & & \multicolumn{2}{p{0.3\linewidth}}{
        \cellcolor[rgb]{0.9,0.9,0.9}{
            \makecell[{{p{\linewidth}}}]{
                \texttt{\tiny{[GM$|$GM]}}
                \texttt{Maximum number of rounds reached} \\
            }
        }
    }
    & & \\ \\

    \theutterance \stepcounter{utterance}  
    & & & \multicolumn{2}{p{0.3\linewidth}}{
        \cellcolor[rgb]{0.9,0.9,0.9}{
            \makecell[{{p{\linewidth}}}]{
                \texttt{\tiny{[GM$|$GM]}}
                \texttt{Initial states:} \\
\texttt{```} \\
\texttt{Player 1:} \\
\texttt{ 1234567} \\
\texttt{╔══╤════╗} \\
\texttt{║◌◌│C◌◌◌║ 1} \\
\texttt{╟──┤◌L┌─╢ 2} \\
\texttt{║◌◌├──┤◌║ 3} \\
\texttt{╟──┘◌◌├─╢ 4} \\
\texttt{║◌◌P◌◌│◌║ 5} \\
\texttt{╟─┐◌┌─┼─╢ 6} \\
\texttt{║◌│◌│◌│◌║ 7} \\
\texttt{╚═╧═╧═╧═╝} \\
\\ 
\\ 
\texttt{Player 2:} \\
\texttt{ 1234567} \\
\texttt{╔══╤════╗} \\
\texttt{║◌◌│L◌◌C║ 1} \\
\texttt{╟──┤◌◌┌─╢ 2} \\
\texttt{║◌◌├──┤◌║ 3} \\
\texttt{╟──┘◌◌├─╢ 4} \\
\texttt{║◌◌P◌◌│◌║ 5} \\
\texttt{╟─┐◌┌─┼─╢ 6} \\
\texttt{║◌│◌│◌│◌║ 7} \\
\texttt{╚═╧═╧═╧═╝} \\
\\ 
\texttt{```} \\
            }
        }
    }
    & & \\ \\

    \theutterance \stepcounter{utterance}  
    & & & \multicolumn{2}{p{0.3\linewidth}}{
        \cellcolor[rgb]{0.9,0.9,0.9}{
            \makecell[{{p{\linewidth}}}]{
                \texttt{\tiny{[GM$|$GM]}}
                \texttt{End states:} \\
\texttt{```} \\
\texttt{Player 1:} \\
\texttt{ 1234567} \\
\texttt{╔══╤════╗} \\
\texttt{║◌◌│C◌L◌║ 1} \\
\texttt{╟──┤◌◌┌─╢ 2} \\
\texttt{║◌◌├──┤◌║ 3} \\
\texttt{╟──┘◌◌├─╢ 4} \\
\texttt{║◌◌P◌◌│◌║ 5} \\
\texttt{╟─┐◌┌─┼─╢ 6} \\
\texttt{║◌│◌│◌│◌║ 7} \\
\texttt{╚═╧═╧═╧═╝} \\
\\ 
\\ 
\texttt{Player 2:} \\
\texttt{ 1234567} \\
\texttt{╔══╤════╗} \\
\texttt{║C◌│◌◌L◌║ 1} \\
\texttt{╟──┤◌◌┌─╢ 2} \\
\texttt{║◌◌├──┤◌║ 3} \\
\texttt{╟──┘◌◌├─╢ 4} \\
\texttt{║◌◌P◌◌│◌║ 5} \\
\texttt{╟─┐◌┌─┼─╢ 6} \\
\texttt{║◌│◌│◌│◌║ 7} \\
\texttt{╚═╧═╧═╧═╝} \\
\\ 
\texttt{```} \\
            }
        }
    }
    & & \\ \\

    \theutterance \stepcounter{utterance}  
    & & & \multicolumn{2}{p{0.3\linewidth}}{
        \cellcolor[rgb]{0.9,0.9,0.9}{
            \makecell[{{p{\linewidth}}}]{
                \texttt{\tiny{[GM$|$GM]}}
                \texttt{* success: True} \\
\texttt{* lose: False} \\
\texttt{* aborted: False} \\
\texttt{{-}{-}{-}{-}{-}{-}{-}} \\
\texttt{* Shifts: 1.00} \\
\texttt{* Max Shifts: 4.00} \\
\texttt{* Min Shifts: 2.00} \\
\texttt{* End Distance Sum: 3.00} \\
\texttt{* Init Distance Sum: 4.41} \\
\texttt{* Expected Distance Sum: 12.57} \\
\texttt{* Penalties: 5.00} \\
\texttt{* Max Penalties: 8.00} \\
\texttt{* Rounds: 12.00} \\
\texttt{* Max Rounds: 12.00} \\
\texttt{* Object Count: 3.00} \\
            }
        }
    }
    & & \\ \\

\end{supertabular}
}

\end{document}
